\section{Introduction}

A public count package has two different integrity problems. First, the package
must preserve the meaning of public election records: contests, reporting units,
batches, statuses, residual counts, and source references. Second, it must make
later byte-level changes visible. A normalized summary can still add correctly
after the raw source file has been edited, replaced, truncated, or omitted.

RCOUNT therefore separates arithmetic checks from hash checks. A verifier may
report that local contest sums and jurisdiction totals pass while the source
hash check fails. That is not a contradiction. It means the normalized records
reconcile, but the raw evidence no longer matches the package's declared source
index.

The discipline is tamper evidence, not omniscience. Hashes can show that bytes
changed. They cannot show that the original bytes were lawful, complete,
accurately exported, free of malware, or sufficient for certification.
